\chapter{Vergleichen, Absorbieren und Ersetzen}
\label{b14guide:ch:calculus}

\section{Vergleich zweier Paarwerte}

\begin{theorem}[Monotonie in beiden Operanden]
\label{b14guide:thm:monotone}
Seien $x\leq x'$ und $y\leq y'$. Sind die beiden
Paarinfima vorhanden, so gilt
\[
 x\wedge y\leq x'\wedge y'.
\]
Sind die beiden Paarsuprema vorhanden, so gilt
\[
 x\vee y\leq x'\vee y'.
\]
Insbesondere gilt dies bei Änderung nur eines Operanden.
Die Existenz eines der verglichenen Werte allein genügt nicht.
\end{theorem}

\begin{theorem}[Vergleich und Gleichheit durch Operanden]
\label{b14guide:thm:compare-pairs}
Sind $x\wedge y$ und $u\wedge v$ vorhanden, so gilt
\[
 x\wedge y\leq u\wedge v
 \Longleftrightarrow
 (x\wedge y\leq u\land x\wedge y\leq v).
\]
Sind $x\vee y$ und $u\vee v$ vorhanden, so gilt
\[
 x\vee y\leq u\vee v
 \Longleftrightarrow
 (x\leq u\vee v\land y\leq u\vee v).
\]
Die jeweiligen Werte sind genau dann gleich, wenn das
angegebene Kriterium in beiden Richtungen erfüllt ist.
\end{theorem}

\begin{theorem}[Ein Paarsupremum unter einem Paarinfimum]
\label{b14guide:thm:cross}
Sind $x\vee y$ und $u\wedge v$ vorhanden, so gilt
\[
 x\vee y\leq u\wedge v
 \Longleftrightarrow
 (x\leq u\land x\leq v\land y\leq u\land y\leq v).
\]
Sind $x\wedge y$ und $x\vee z$ vorhanden, so gilt stets
\[
 x\wedge y\leq x\vee z.
\]
\end{theorem}

\section{Absorptionsregeln}

\begin{theorem}[Absorption mit nur einem vorgegebenen Paarwert]
\label{b14guide:thm:absorption}
Ist $x\vee y$ vorhanden, so existieren auch die folgenden
Paarinfima und haben die Werte
\[
 x\wedge(x\vee y)=x,\qquad
 y\wedge(x\vee y)=y.
\]
Ist $x\wedge y$ vorhanden, so existieren auch die folgenden
Paarsuprema und haben die Werte
\[
 x\vee(x\wedge y)=x,\qquad
 y\vee(x\wedge y)=y.
\]
Die Existenz des jeweils anderen ursprünglichen Paarwertes
wird nicht vorausgesetzt.
\end{theorem}

\begin{theorem}[Wiederholte Operanden]
\label{b14guide:thm:repeated}
Ist $m=x\wedge y$ vorhanden, so existieren
\[
 x\wedge m=m,\qquad y\wedge m=m.
\]
Ist $s=x\vee y$ vorhanden, so existieren
\[
 x\vee s=s,\qquad y\vee s=s.
\]
Zusammen mit der Kommutativität ergeben sich die entsprechenden
Formeln bei vertauschter Stellung der Operanden.
\end{theorem}

\section{Ersetzen eines Operanden}

\begin{theorem}[Ersetzen zwischen Operand und Paarwert]
\label{b14guide:thm:replacement}
Ist $m=x\wedge y$ vorhanden und $m\leq x'\leq x$, so existiert
\[
 x'\wedge y=m.
\]
Ist $s=x\vee y$ vorhanden und $x\leq x'\leq s$, so existiert
\[
 x'\vee y=s.
\]
Entsprechend darf der zweite Operand innerhalb der genannten
Grenzen ersetzt werden.
\end{theorem}

\begin{theorem}[Ersetzen zwischen zwei gleichen Paarwerten]
\label{b14guide:thm:plateau}
Seien $x\leq z\leq y$. Existieren $x\wedge a$ und $y\wedge a$
und sind sie gleich $m$, so existiert auch $z\wedge a=m$.
Existieren $x\vee a$ und $y\vee a$ und sind sie gleich $s$,
so existiert auch $z\vee a=s$.
\end{theorem}